% GENERATED FILE. Edit canonical lessons, then run npm run generate:books.
% canonical-source-hash: fnv1a64:19d71fd8ecaa4361
% canonical-lessons: KA-C348-nidhi, KA-C348-kadalgalla, KA-C348-sahasa, KA-C348-rahasya, KA-C348-mantrika

\chapter{The Pirate Treasure}
\label{ch:ka-the-pirate-treasure}
\hlchaptermodality{348}

\begin{chapteropening}
\textbf{By the end of this chapter:} \emph{I can say a treasure, a pirate, an adventure, a secret and a sorcerer.}

Complete the last lesson of chapter 348: I can say a treasure, a pirate, an adventure, a secret and a sorcerer.
\end{chapteropening}

\section[\texorpdfstring{\kn{ನಿಧಿ} (nidhi)}{ನಿಧಿ (nidhi)}]{\kn{ನಿಧಿ} (nidhi) — a treasure}

\label{lesson:KA-C348-nidhi}



\begin{quote}
Before the new one: say the Kannada for to mention, then the Kannada for a summary.
\end{quote}

\subsection*{You'll want to know: \kn{ನಿಧಿ}}
\textbf{\kn{ನಿಧಿ}} — \emph{nidhi} — \textquotedblleft{}a treasure\textquotedblright{}.

\begin{quote}
\kn{ಕಥೆಯಲ್ಲಿ} \kn{ಹುಡುಗರಿಗೆ} \kn{ಗುಹೆಯಲ್ಲಿ} \kn{ನಿಧಿ} \kn{ಸಿಕ್ಕಿತು}.
\end{quote}

\begin{quote}
\emph{katheyalli huḍugarige guheyalli nidhi sikkitu.}
\end{quote}

\begin{quote}
In the story the boys found a treasure in a cave.
\end{quote}

Plural: \kn{ನಿಧಿಗಳು}, \emph{nidhigaḷu}.

Also a fund: \kn{ಕಲ್ಯಾಣ} \kn{ನಿಧಿ}, a welfare fund.

\begin{cousinweb}
\textbf{\kn{ನಿಧಿ}} (\emph{nidhi}) is borrowed from Sanskrit \emph{nidhi}, \textquotedblleft{}setting down, a hoard\textquotedblright{} (CDIAL 7207), from \emph{ni-dhā}, \textquotedblleft{}to put down\textquotedblright{}.
\end{cousinweb}

\begin{grammarlens}[title={What you've built}]
One more word for telling what happened.
\end{grammarlens}

\subsection*{Your turn}
\begin{itemize}
\raggedright
  \item \emph{Say it:} \emph{nidhi}
  \item \emph{Say it:} \emph{nidhi}, once more
  \item \emph{Say it:} \emph{katheyalli huḍugarige guheyalli nidhi sikkitu.} — In the story the boys found a treasure in a cave.
  \item \emph{Recall:} say the Kannada for to mention, then the Kannada for a summary, then say \emph{nidhi} again
\end{itemize}

\begin{tcolorbox}[breakable,colback=teal!4,colframe=teal!35!black,title={Before you move on}]
What does \kn{ನಿಧಿ} mean? (\textquotedblleft{}A treasure\textquotedblright{}.) Say it once more.
\end{tcolorbox}
\section[\texorpdfstring{\kn{ಕಡಲ್ಗಳ್ಳ} (kaḍalgaḷḷa)}{ಕಡಲ್ಗಳ್ಳ (kaḍalgaḷḷa)}]{\kn{ಕಡಲ್ಗಳ್ಳ} (kaḍalgaḷḷa) — a pirate}

\label{lesson:KA-C348-kadalgalla}



\begin{quote}
Before the new one: say the Kannada for a summary, then the Kannada for a treasure.
\end{quote}

\subsection*{You'll want to know: \kn{ಕಡಲ್ಗಳ್ಳ}}
\textbf{\kn{ಕಡಲ್ಗಳ್ಳ}} — \emph{kaḍalgaḷḷa} — \textquotedblleft{}a pirate\textquotedblright{}.

\begin{quote}
\kn{ಕಡಲ್ಗಳ್ಳರು} \kn{ಹಡಗಿನ} \kn{ಚಿನ್ನವನ್ನು} \kn{ದೋಚಿದರು}.
\end{quote}

\begin{quote}
\emph{kaḍalgaḷḷaru haḍagina cinnavannu dōcidaru.}
\end{quote}

\begin{quote}
The pirates looted the ship's gold.
\end{quote}

Plural: \kn{ಕಡಲ್ಗಳ್ಳರು}, \emph{kaḍalgaḷḷaru}.

\kn{ಕಡಲ್}, sea, and \kn{ಕಳ್ಳ}, thief: a thief of the sea.

\begin{cousinweb}
\textbf{\kn{ಕಡಲ್ಗಳ್ಳ}} (\emph{kaḍalgaḷḷa}) is a native Dravidian compound: \emph{kaḍal}, \textquotedblleft{}the sea\textquotedblright{} (DEDR 1118), and \emph{kaḷḷa}, \textquotedblleft{}a thief\textquotedblright{} (DEDR 1372).
\end{cousinweb}

\begin{grammarlens}[title={What you've built}]
One more word for telling what happened.
\end{grammarlens}

\subsection*{Your turn}
\begin{itemize}
\raggedright
  \item \emph{Say it:} \emph{kaḍalgaḷḷa}
  \item \emph{Say it:} \emph{kaḍalgaḷḷa}, once more
  \item \emph{Say it:} \emph{kaḍalgaḷḷaru haḍagina cinnavannu dōcidaru.} — The pirates looted the ship's gold.
  \item \emph{Recall:} say the Kannada for a summary, then the Kannada for a treasure, then say \emph{kaḍalgaḷḷa} again
\end{itemize}

\begin{tcolorbox}[breakable,colback=teal!4,colframe=teal!35!black,title={Before you move on}]
What does \kn{ಕಡಲ್ಗಳ್ಳ} mean? (\textquotedblleft{}A pirate\textquotedblright{}.) Say it once more.
\end{tcolorbox}
\section[\texorpdfstring{\kn{ಸಾಹಸ} (sāhasa)}{ಸಾಹಸ (sāhasa)}]{\kn{ಸಾಹಸ} (sāhasa) — an adventure; a daring act}

\label{lesson:KA-C348-sahasa}



\begin{quote}
Before the new one: say the Kannada for a treasure, then the Kannada for a pirate.
\end{quote}

\subsection*{You'll want to know: \kn{ಸಾಹಸ}}
\textbf{\kn{ಸಾಹಸ}} — \emph{sāhasa} — \textquotedblleft{}an adventure; a daring act\textquotedblright{}.

\begin{quote}
\kn{ಕಾಡಿನಲ್ಲಿ} \kn{ನಮ್ಮ} \kn{ಮೂರು} \kn{ದಿನದ} \kn{ಸಾಹಸ} \kn{ಮರೆಯಲು} \kn{ಆಗುವುದಿಲ್ಲ}.
\end{quote}

\begin{quote}
\emph{kāḍinalli namma mūru dinada sāhasa mareyalu āguvudilla.}
\end{quote}

\begin{quote}
Our three-day adventure in the forest is impossible to forget.
\end{quote}

Plural: \kn{ಸಾಹಸಗಳು}, \emph{sāhasagaḷu}.

A daring person is a \kn{ಸಾಹಸಿ}.

\begin{cousinweb}
\textbf{\kn{ಸಾಹಸ}} (\emph{sāhasa}) is borrowed from Sanskrit \emph{sāhasa}, \textquotedblleft{}boldness, rashness\textquotedblright{} (CDIAL 13378).
\end{cousinweb}

\begin{grammarlens}[title={What you've built}]
One more word for telling what happened.
\end{grammarlens}

\subsection*{Your turn}
\begin{itemize}
\raggedright
  \item \emph{Say it:} \emph{sāhasa}
  \item \emph{Say it:} \emph{sāhasa}, once more
  \item \emph{Say it:} \emph{kāḍinalli namma mūru dinada sāhasa mareyalu āguvudilla.} — Our three-day adventure in the forest is impossible to forget.
  \item \emph{Recall:} say the Kannada for a treasure, then the Kannada for a pirate, then say \emph{sāhasa} again
\end{itemize}

\begin{tcolorbox}[breakable,colback=teal!4,colframe=teal!35!black,title={Before you move on}]
What does \kn{ಸಾಹಸ} mean? (\textquotedblleft{}An adventure; a daring act\textquotedblright{}.) Say it once more.
\end{tcolorbox}
\section[\texorpdfstring{\kn{ರಹಸ್ಯ} (rahasya)}{ರಹಸ್ಯ (rahasya)}]{\kn{ರಹಸ್ಯ} (rahasya) — a secret; a mystery}

\label{lesson:KA-C348-rahasya}



\begin{quote}
Before the new one: say the Kannada for a pirate, then the Kannada for an adventure.
\end{quote}

\subsection*{You'll want to know: \kn{ರಹಸ್ಯ}}
\textbf{\kn{ರಹಸ್ಯ}} — \emph{rahasya} — \textquotedblleft{}a secret; a mystery\textquotedblright{}.

\begin{quote}
\kn{ನಿಧಿಯ} \kn{ರಹಸ್ಯ} \kn{ಒಂದು} \kn{ಹಳೆಯ} \kn{ನಕ್ಷೆಯಲ್ಲಿ} \kn{ಇತ್ತು}.
\end{quote}

\begin{quote}
\emph{nidhiya rahasya ondu haḷeya nakṣeyalli ittu.}
\end{quote}

\begin{quote}
The treasure's secret was in an old map.
\end{quote}

Plural: \kn{ರಹಸ್ಯಗಳು}, \emph{rahasyagaḷu}.

As an adjective: \kn{ರಹಸ್ಯ} \kn{ದಾರಿ}, a secret path. \kn{ರಹಸ್ಯವಾಗಿ}, secretly.

\begin{cousinweb}
\textbf{\kn{ರಹಸ್ಯ}} (\emph{rahasya}) is borrowed from Sanskrit \emph{rahasya}, \textquotedblleft{}private, secret\textquotedblright{} (CDIAL 10669), from \emph{rahas}, \textquotedblleft{}a lonely spot\textquotedblright{}.
\end{cousinweb}

\begin{grammarlens}[title={What you've built}]
One more word for telling what happened.
\end{grammarlens}

\subsection*{Your turn}
\begin{itemize}
\raggedright
  \item \emph{Say it:} \emph{rahasya}
  \item \emph{Say it:} \emph{rahasya}, once more
  \item \emph{Say it:} \emph{nidhiya rahasya ondu haḷeya nakṣeyalli ittu.} — The treasure's secret was in an old map.
  \item \emph{Recall:} say the Kannada for a pirate, then the Kannada for an adventure, then say \emph{rahasya} again
\end{itemize}

\begin{tcolorbox}[breakable,colback=teal!4,colframe=teal!35!black,title={Before you move on}]
What does \kn{ರಹಸ್ಯ} mean? (\textquotedblleft{}A secret; a mystery\textquotedblright{}.) Say it once more.
\end{tcolorbox}
\section[\texorpdfstring{\kn{ಮಾಂತ್ರಿಕ} (māntrika)}{ಮಾಂತ್ರಿಕ (māntrika)}]{\kn{ಮಾಂತ್ರಿಕ} (māntrika) — a sorcerer; a magician}

\label{lesson:KA-C348-mantrika}



\begin{quote}
Before the new one: say the Kannada for an adventure, then the Kannada for a secret.
\end{quote}

\subsection*{You'll want to know: \kn{ಮಾಂತ್ರಿಕ}}
\textbf{\kn{ಮಾಂತ್ರಿಕ}} — \emph{māntrika} — \textquotedblleft{}a sorcerer; a magician\textquotedblright{}.

\begin{quote}
\kn{ಮಾಂತ್ರಿಕ} \kn{ಆ} \kn{ನಿಧಿಯನ್ನು} \kn{ಕಾಯುತ್ತಿದ್ದನು}.
\end{quote}

\begin{quote}
\emph{māntrika ā nidhiyannu kāyuttiddanu.}
\end{quote}

\begin{quote}
A sorcerer was guarding that treasure.
\end{quote}

Plural: \kn{ಮಾಂತ್ರಿಕರು}, \emph{māntrikaru}.

One who works spells, \kn{ಮಂತ್ರ}, in a story.

\begin{cousinweb}
\textbf{\kn{ಮಾಂತ್ರಿಕ}} (\emph{māntrika}) is borrowed from Sanskrit \emph{māntrika}, \textquotedblleft{}a reciter of spells, a sorcerer\textquotedblright{} (CDIAL 10050), from \emph{mantra}, \textquotedblleft{}a spell\textquotedblright{}.
\end{cousinweb}

\begin{grammarlens}[title={What you've built}]
One more word for telling what happened.
\end{grammarlens}

\subsection*{Your turn}
\begin{itemize}
\raggedright
  \item \emph{Say it:} \emph{māntrika}
  \item \emph{Say it:} \emph{māntrika}, once more
  \item \emph{Say it:} \emph{māntrika ā nidhiyannu kāyuttiddanu.} — A sorcerer was guarding that treasure.
  \item \emph{Recall:} say the Kannada for an adventure, then the Kannada for a secret, then say \emph{māntrika} again
\end{itemize}

\begin{tcolorbox}[breakable,colback=teal!4,colframe=teal!35!black,title={Before you move on}]
What does \kn{ಮಾಂತ್ರಿಕ} mean? (\textquotedblleft{}A sorcerer; a magician\textquotedblright{}.) Say it once more.
\end{tcolorbox}
